\documentclass[exercice,a4]{thedocument}%
\usepackage[french]{babel}%
\usepackage{scratch3}%
\startdatakeys
    \source{21GENMATMEAG1}
    \BO{2026}
    \cycle{4}
    \competencesDuSocle{CA,RA}
    \objectifsApprentissage{C4D2ba,C4A2ed}
    \motsCles{programme de calcul,Scratch,expression littérale,factorisation,équation produit}
\enddatakeys
\begin{document}%

\medskip

\parbox{0.45\linewidth}{Voici un programme de calcul:}
\hfill
\parbox{0.45\linewidth}{\begin{tabular}{|l|}\hline
Choisir un nombre.\\
Prendre le carré du nombre de départ.\\
Ajouter le triple du nombre de départ.\\
Soustraire 10 au résultat.\\ \hline
\end{tabular}}

\medskip

\begin{enumerate}
\item Vérifier que si on choisit 4 comme nombre de départ, on obtient 18.
\item Appliquer ce programme de calcul au nombre $- 3$.
\item Vous trouverez ci-dessous un script, écrit avec scratch.

\begin{scratch}[num blocks]
\blockinit{Quand \greenflag est cliqué}
\blocksensing{demander \ovalnum{Choisis un nombre } et attendre}
\blockpen{mettre \selectmenu{x} à \ovalsensing{Réponse}}
\blockpen{mettre \selectmenu{y} à \ovaloperator{\ovalnum{x}*\ovalnum{x}}}
\blockpen{mettre \selectmenu{z} à \ovalnum{y} + \ovaloperator{\ovalnum{}*\ovalnum{}}}
\blockpen{mettre \selectmenu{Résultat} à \ovaloperator{\ovalnum{} - \ovalnum{}}}
\blocklook{dire \ovaloperator{regroupe \ovalnum{Le nombre final}  est \ovaloperator{Résultat} pendant \ovalnum{2} secondes}}
\end{scratch}

\medskip

Compléter sur l'ANNEXE  les lignes 5 et 6 pour que ce script corresponde au programme de calcul.

\item On veut déterminer le nombre à choisir au départ pour obtenir zéro comme résultat.
	\begin{enumerate}
		\item On appelle $x$ le nombre de départ. Exprimer en fonction de $x$ le résultat final.
		\item Vérifier que ce résultat peut aussi s'écrire sous la forme $(x + 5)(x - 2)$.
		\item Quel(s) nombre(s) doit-on choisir au départ pour obtenir le nombre $0$ à l'arrivée ?
	\end{enumerate}
\end{enumerate}

\bigskip

\vspace{2cm}

\begin{scratch}[num blocks]
\blockinit{Quand \greenflag est cliqué}
\blocksensing{demander \ovalnum{Choisis un nombre } et attendre}
\blockpen{mettre \selectmenu{x} à \ovalsensing{Réponse}}
\blockpen{mettre \selectmenu{y} à \ovaloperator{\ovalnum{x}*\ovalnum{x}}}
\blockpen{mettre \selectmenu{z} à \ovalnum{y} + \ovaloperator{\ovalnum{}*\ovalnum{}}}
\blockpen{mettre \selectmenu{Résultat} à \ovaloperator{\ovalnum{} - \ovalnum{}}}
\blocklook{dire \ovaloperator{regroupe \ovalnum{Le nombre final} est \ovaloperator{Résultat} pendant \ovalnum{2} secondes}}
\end{scratch}

\showthesource
\end{document}
